\section{*Electrons in a uniform electric field (12/2)}


\subsection{Solving Schrodinger equation in a uniform electric field}

Let the electric field be $\ce\hat{\bm x}$ and define the electric force
\begin{equation}
 F=e\ce.
\end{equation}
for an electron of charge $e$. In one dimension, the electric field $\ce$ is related to scalar and vector potentials by
\begin{align}
    \ce=-\nabla\phi-\partial_t A
\end{align}

\subsubsection{Gauge 1: time-dependent vector-potential gauge}
In one dimension,
\begin{equation}
\hat H=\frac{[-\imth\hbar\partial_x-eA(x,t)]^2}{2m}+e\phi(x,t).
\end{equation}
Choose the following gauge:
\begin{equation}
\phi_1=0,
\qquad
A_1=-\mathcal Et.
\end{equation}
Then
\begin{equation}
\imth\hbar\frac{\dif}{\dif t}\dket\psi
=\frac{(\hat p+e\mathcal Et)^2}{2m}\dket\psi.
\end{equation}
For initial state
\begin{equation}
\psi(x,0)=e^{\imth kx},
\end{equation}
one obtains
\begin{equation}
\boxed{
\psi(x,t)=
\exp\left[
\imth kx-\frac{\imth}{6me\mathcal E\hbar}
\left((\hbar k+e\mathcal Et)^3-(\hbar k)^3\right)
\right].}
\end{equation}

\subsubsection{Gauge 2: scalar-potential gauge and Airy functions}
Alternatively we can choose
\begin{equation}
\phi_2=-\mathcal Ex,
\qquad
A_2=0.
\end{equation}
Then
\begin{equation}
\hat H=\frac{\hat p^2}{2m}-e\mathcal E\hat x.
\end{equation}
The stationary equation is
\begin{equation}
-\frac{\hbar^2}{2m}\psi''(x)-e\mathcal Ex\psi(x)=E\psi(x).
\end{equation}
With
\begin{equation}
u=\left(\frac{2m|e|\mathcal E}{\hbar^2}\right)^{1/3}
\left(x-\frac{E}{e\mathcal E}\right),
\end{equation}
one obtains the Airy equation
\begin{equation}
\frac{\dif^2\psi}{\dif u^2}=u\psi.
\end{equation}
The decaying solution is
\begin{equation}
\boxed{
\operatorname{Ai}(u)=\frac1\pi\int_0^\infty
\cos\left(\frac{t^3}{3}+ut\right)\dif t.}
\end{equation}
Its asymptotic forms are
\begin{align}
\operatorname{Ai}(u)&\sim
\frac{1}{2\sqrt\pi}u^{-1/4}e^{-2u^{3/2}/3},
&&u\gg1,\\
\operatorname{Ai}(u)&\sim
\frac{1}{\sqrt\pi}|u|^{-1/4}
\cos\left(\frac23|u|^{3/2}-\frac\pi4\right),
&&u\ll-1.
\end{align}
The spectrum is continuous.

\subsubsection{Relation between the two gauges}
Use the gauge convention
\begin{equation}
 A'=A-\partial_xf,
 \qquad \phi'=\phi+\partial_tf,
 \qquad \psi'=e^{-\imth qf/\hbar}\psi.
\end{equation}
The choice
\begin{equation}
 \boxed{f(x,t)=-\ce tx}
\end{equation}
transforms Gauge 1 into Gauge 2:
\begin{equation}
 A_2=A_1-\partial_xf=0,
 \qquad
 \phi_2=\phi_1+\partial_tf=-\ce x.
\end{equation}
Therefore,
\begin{equation}
 \boxed{\psi_2(x,t)=e^{\imth Ftx/\hbar}\psi_1(x,t).}
\end{equation}
The Hamiltonians satisfy
\begin{equation}
 \hat H_2=\hat U_f\hat H_1\hat U_f^\dagger
 +\imth\hbar\dot{\hat U}_f\hat U_f^\dagger,
 \qquad
 \hat U_f=e^{\imth Ft\hat x/\hbar}.
\end{equation}
The transformed plane-wave solution in gauge 2 is
\begin{equation}
 \boxed{
 \psi_2(x,t)=C
 \exp\left[
 \imth\left(k+\frac{Ft}{\hbar}\right)x
 -\frac{\imth}{2m\hbar}\int_0^t(\hbar k+F\tau)^2\dif\tau
 \right].}
\end{equation}
It is a superposition of stationary Airy eigenstates $\{\dket{\varphi_E}\}$:
\begin{equation}
 \psi_2(x,t)=\int_{-\infty}^{\infty}
 c(E)e^{-\imth Et/\hbar}\varphi_E(x)\dif E.
\end{equation}

\subsubsection{Density and electric current in the two gauges}
For the plane-wave solution in Gauge 1, the charge density and electric currents are
\begin{align}
% \rho_1&=|C|^2,\\
 \rho_{q,1}&=e|C|^2,\\
% j_1&=\frac{|C|^2}{m}(\hbar k+Ft),\\
 j_{q,1}&=\frac{e|C|^2}{m}(\hbar k+Ft).
\end{align}
For the gauge-transformed solution in Gauge 2, we also have
\begin{align}
% \rho_2&=|C|^2,\\
 \rho_{q,2}&=e|C|^2,\\
% j_2&=\frac{|C|^2}{m}(\hbar k+Ft),\\
 j_{q,2}&=\frac{e|C|^2}{m}(\hbar k+Ft).
\end{align}
Hence
\begin{equation}
 \boxed{\rho_{q,1}=\rho_{q,2},
 \qquad j_{q,1}=j_{q,2}.}
\end{equation}

A real stationary Airy state $\dket{\varphi_E}$ in Gauge 2 is
\begin{equation}
 \psi_{2,E}(x,t)=e^{-\imth Et/\hbar}\varphi_E(x),
 \qquad \varphi_E(x)\in\mathbb R.
\end{equation}
Its charge density and current are
\begin{equation}
 \rho_{q,2}^{(E)}=e|\varphi_E(x)|^2,
 \qquad
 \boxed{j_{q,2}^{(E)}=0.}
\end{equation}
The corresponding Gauge-1 state is
\begin{equation}
 \boxed{
 \psi_{1,E}(x,t)=e^{-\imth Ftx/\hbar}
 e^{-\imth Et/\hbar}\varphi_E(x).}
\end{equation}
For this state,
\begin{equation}
 \rho_{q,1}^{(E)}=e|\varphi_E(x)|^2,
 \qquad
 \boxed{j_{q,1}^{(E)}=0.}
\end{equation}
The phase gradient contributes canonical momentum $-Ft$, while the vector
potential contributes $+Ft$ to the kinetic momentum, so the physical current
vanishes in both gauges.

\subsection{Landau--Zener Tunneling in solids}

\subsubsection{Band crossing}
Consider
\begin{equation}
\hat H=\frac{\hat p^2}{2m}+V(\hat x)+|e|\mathcal E\hat x,
\qquad
V(x+a)=V(x).
\end{equation}
Near a valence-conduction avoided crossing, the conduction and valence bands can be modeled by a two-level system, where the energy of conduction band $\dket{c}$ is higher than the valence band $\dket{v}$ by the band gap $E_g>0$:
\begin{equation}
\hat H\simeq
\begin{pmatrix}
E_g/2+\hbar^2k^2/(2m_v)&\Delta\\
\Delta^*&-E_g/2-\hbar^2k^2/(2m_c)
\end{pmatrix}.
\end{equation}
$\Delta$ models the tunneling term between the two bands. 

The electric field gives
\begin{equation}
\hbar\frac{\dif k}{\dif t}\simeq e\mathcal E,
\qquad
k(t)=k_0+\frac{e\mathcal E}{\hbar}t.
\end{equation}

\subsubsection{Landau--Zener model}
In a semiconductor, we expand the wavevector $k$ around the smallest direct gap location $k_0$, and we arrive at the Landau-Zener model
\begin{equation}
\hat H_{\mathrm{LZ}}(t)=
\begin{pmatrix}
\alpha t&\Delta\\
\Delta&-\alpha t
\end{pmatrix},
\qquad
\dket{\psi(t)}=
\begin{pmatrix}C_1(t)\\C_2(t)\end{pmatrix}.
\end{equation}
where
\begin{align}
    \alpha=|e|\mathcal E v_F/\hbar,~~~v_F=\frac{\hbar k_0}{m^\ast},~~~\frac1{m^*}=\frac1{m_c}+\frac1{m_v}.
\end{align}

The 2-component Schrodinger equation is
\begin{align}
\imth\hbar\dot C_1&=\alpha tC_1+\Delta C_2,\\
\imth\hbar\dot C_2&=\Delta C_1-\alpha tC_2.
\end{align}
Eliminating $C_1$ gives
\begin{equation}
-\hbar^2\ddot C_2
=\left(\alpha^2t^2+\Delta^2-\imth\hbar\alpha\right)C_2.
\end{equation}






\subsubsection{Computing the tunneling probability}
Define
\begin{equation}
 z=e^{-\imth\pi/4}\sqrt{\frac{2\alpha}{\hbar}}\,t,
 \qquad
 \delta=\frac{\Delta^2}{2\hbar\alpha}.
\end{equation}
Then
\begin{equation}
 \boxed{
 \frac{\dif^2C_2}{\dif z^2}
 +\left(\nu+\frac12-\frac{z^2}{4}\right)C_2=0,
 \qquad \nu=\imth\delta.}
\end{equation}
This is Weber's equation. Its solutions are parabolic-cylinder functions $D_\nu(z)$,
also called Weber functions:
\begin{equation}
 \boxed{C_2(z)=A D_{\imth\delta}(z)
 +B D_{\imth\delta}(-z).}
\end{equation}
The remaining amplitude follows from
\begin{equation}
 \boxed{C_1(t)=\frac{\imth\hbar\dot C_2+\alpha tC_2}{\Delta}.}
\end{equation}

If the initial state is $\dket{\psi(t=-\infty)}=(0,1)^T$
\begin{equation}
C_2(-\infty)=1,
\qquad
C_1(-\infty)=0,
\end{equation}
then
\begin{equation}
|C_2(+\infty)|^2
=e^{-2\pi\Delta^2/(\hbar\alpha)}\equiv P_\text{diabatic}.
\end{equation}
Hence the transition probability is
\begin{equation}
\boxed{
P_{1\to2}=1-e^{-2\pi\Delta^2/(\hbar\alpha)}.}
\end{equation}
The diabatic probability (Zener excitation probability) $P_\text{diabatic}$ increases, as parameter $\alpha$ increases. 

In other words, as the electric field $\ce$ (and hence $\alpha$) increases, there is a larger chance for electrons to diabatically transfer from the lower-energy valence band to higher-energy conduction band. 


